\section{The fifteen assertions in the automorphy literature hypothesis}

% BEGIN A72F1
\paragraph{A72F1. The two Artin maps and the ray-principal test.}
Let $E$ be a number field and take the specified algebraic closure.
CIP1 gives the five clauses of $\mathrm{IsGlobalArtin}$ with the
geometric convention. In particular a uniformizer maps to the coset of
inverse arithmetic Frobenius, local units map onto the image of inertia,
and the closure of the local image gives the image of the decomposition
group. This supplies $\mathrm{art}$.

The global map in this description is unique. Two such maps have
the same kernel and are quotient maps, so they differ by a continuous
automorphism $b$ of $G_E^{\rm ab}$. For each finite quotient $q$,
both $q$ and $qb$ are unramified outside finitely many places. At a
place unramified for both, the uniformizer clause gives
$q(F_v)=qb(F_v)$ for geometric Frobenius. Apply Chebotarev to
the finite quotient by $\ker q\cap\ker(qb)$. This proves $q=qb$
for every $q$ and hence $b=1$. For every admissible local map
$a:W_w\to E_w^\times$, its global compatibility therefore identifies
its composite with the global map on the local idele with the given
abelianized Weil element. This identity suffices for rank-one
characters below.

For a nonzero finite ideal $\mathfrak m$, take the ray class field
with modulus $\mathfrak m$ times all real places. Finite Artin
reciprocity gives the ray class group isomorphism, sending
a prime class to \emph{arithmetic} Frobenius as stipulated in
the separate $\mathrm{ray}$ field. The primes away from
$\mathfrak m$ are unramified and their Frobenius elements generate
the Galois group. Its finite complex characters separate points.
Thus an integral ideal $I$ prime to $\mathfrak m$ has value one
under all these characters precisely when $I=(x)$ with $x$ totally
positive and congruent to $1$ at the prescribed powers of all
primes dividing $\mathfrak m$. Choose a totally positive denominator
$\beta$ clearing $x$, prime to $\mathfrak m$: the denominator
ideal is prime to $\mathfrak m$, so take an element in it congruent
to one modulo $\mathfrak m$ and add a sufficiently large positive
rational integer belonging to both ideals. Then $\alpha=x\beta$ is
integral, $\alpha\equiv\beta$ modulo $\mathfrak m$, and
$I(\beta)=(\alpha)$. The converse follows by division. This proves
$\mathrm{RayClassFieldSt}$, including the unit ideal modulus and
the signs at real places.
% END A72F1

% BEGIN A72F2
\paragraph{A72F2. Semistable periods on every tied model.}
Fix the source's $P,AI$ and a witness of $\mathrm{LocTied}$.
As in F71A3 its comparison on $B_{\rm cris}$ and $B_{\rm dR}$
preserves the actions, the Frobenius, the rational scalars and the
filtration. The semistable ring here is explicitly $P.B[X]$.
Fix the source's compatible roots $\ell^{1/\ell^m}$ and its
roots of unity $\zeta_m$. The Kummer cocycle $k(g)$ is defined by
$g(\ell^{1/\ell^m})/\ell^{1/\ell^m}=\zeta_m^{k(g)\bmod\ell^m}$.
On the polynomial generator the source operators are
\[
  gX=X+k(g)t,\qquad \varphi X=\ell X,\qquad N=-d/dX .
\]
These are Fontaine's polynomial presentation of $B_{\rm st}$ with
the logarithm of the chosen compatible roots as generator. The
period $t$ uses the same $\zeta_m$ by the $\zeta$ equality in
$AI$ and the logarithm clause of the tie. Consequently extending
the crystalline comparison by this generator intertwines all three
operators. It respects the unramified scalars: the embedded
prime-to-$\ell$ roots of unity are the unique lifts of their
residues, and their finite unramified fields exhaust
$\mathbb Q_\ell^{\rm nr}$. This comparison identifies the union
of the invariant subspaces over finite extensions with the
usual $D_{\rm pst}$. At a finite extension $K$ of $\mathbb Q_\ell$
and for a continuous finite-dimensional rational representation,
Berger's $p$-adic monodromy theorem (\emph{Repr\'esentations
$p$-adiques et \'equations diff\'erentielles}, Invent. Math.
148 (2002)) gives potential semistability from the de Rham
condition. After scalar extension to $\mathbb Q_\ell^{\rm nr}$
the period module has dimension $\dim_{\mathbb Q_\ell}V$,
exactly $\mathrm{IsPst}$. This proves $\mathrm{berger}$.

For coefficient representations the comparisons commute with
extension of a finite coefficient field, restriction of scalars
and projection to an embedding factor. In particular ranks, each
labelled filtration jump, and the Frobenius-semisimplified
Weil--Deligne class computed from $D_{\rm pst}$ have their
usual meanings on every such model. The place comparison is along
the supplied $\mathrm{LocBridge}$: its inducing embedding transports
the decomposition group, and $\mathrm{IsInduced}$ specifies the
embedding on which the filtration is evaluated. All finite comparison
fields can be enlarged to include the coefficient field of that
bridge.
% END A72F2

% BEGIN A72F3
\paragraph{A72F3. The two local conclusions for algebraic Hecke characters.}
Let $\mu$ be the supplied continuous avatar of a Hecke character
$\chi$. Fix any allowed bridge, local Artin map $a$, and local
component $\xi$ at its $\ell$-adic place. For $s\in W_w$ set $x=a(s)$.
Insert the local idele of $x$ into $\mathrm{IsLAdicOf}$.
The compatibility in A72F1 and the trivial archimedean component
of this idele give
\[
 \mu(s)=\iota^{-1}(\xi(x))\prod_{j: E_w\hookrightarrow
 \overline{\mathbb Q}_\ell} j(x)^{n_j},\qquad
 n_j=n(\iota j|_E).
\]
Here $j$ runs over continuous embeddings. The local map is the
inverse geometric reciprocity map. Indeed its global compatibility
determines it: the injection of a single finite idele factor into
the abelianized Galois group follows from Chevalley's congruence
theorem for units, as in the source's local Artin uniqueness
argument. Choose a uniformizer $\varpi$ and write
$x=\varpi^d u$. On geometric reciprocity the Lubin--Tate character
for $\varpi$ takes $x$ to $u$. Its $j$-component is crystalline,
has filtration jump $-1$ at $j$ and zero at the other embeddings,
and its Weil--Deligne character takes a geometric Frobenius to
$j(\varpi)^{-1}$. These are the rank-one comparisons from the
filtered Frobenius module of the Lubin--Tate group (Conrad,
\emph{Lifting global representations with local properties},
Appendix~B). The product of the labelled powers accounts for
every algebraic factor on units. The quotient character is smooth on
inertia, with value
$\iota^{-1}(\xi(\varpi))\prod_j j(\varpi)^{n_j}$
on the element corresponding to $\varpi$.
A rank-one character with finite inertia is potentially unramified,
hence potentially crystalline with zero jumps. Its Weil character
is the original smooth character on the Weil group. Multiplication
cancels the factors $j(\varpi)^{n_j}$ against the Lubin--Tate
Frobenius factors. The resulting WD character is exactly
$\xi\circ a$, including its unramified factor, with $N=0$.
If $\chi$ is unramified at $w$, the smooth factor has trivial
inertia and the product is crystalline over $E_w$ itself.

These calculations on a finite coefficient field give de Rhamness,
potential crystallinity, crystallinity in the unramified case,
and the singleton jump $\{-n(\iota\tau)\}$ for every allowed
induced label. Enlarge the field to contain the images of the
embeddings, the finite inertia values and the single Frobenius
value. The latter belongs to $\overline{\mathbb Q}_\ell$ by
the formula for the supplied avatar. The continuous character then
takes its values in that field by density of the local Weil group.
Transport using A72F2. This gives $\mathrm{chA}$ and $\mathrm{chW}$.
% END A72F3

% BEGIN A72F4
\paragraph{A72F4. Decay, fixed-type finiteness and the cuspidal Hilbert space.}
The group of a nondegenerate Hermitian space in $\mathrm{decay}$ is
a reductive unitary group with anisotropic rational center. The
hypothesis $\mathrm{IsAutFun}$ supplies a common open compact
finite level. The positive-dimensional isotropic subspaces in
$\mathrm{IsCuspFun}$ index the maximal rational parabolics.
In coordinates adapted to such a subspace $L$, the subgroup tested
there acts identically on $L$, $L^\perp/L$ and $V/L^\perp$.
It is the unipotent radical of its stabilizer. Vanishing of these
constant terms gives all proper constant terms by iterated
integration. The height used here is the product of matrix norms
of $g,g^{-1}$, bounded below by one, with a finite product of
nontrivial finite factors. On a fixed level there are finitely many
arithmetic quotients of the real group. Reduction theory covers
them by finitely many Siegel sets. On each set the finite
representative is fixed up to a compact subgroup and the height
is bounded by a power of a faithful real matrix norm. Conversely
the archimedean norms in the uniform moderate growth assumption
are bounded by a power of that norm. The smooth cuspidal rapid
decay theorem (Miller--Schmid, \emph{On the rapid decay of cuspidal
automorphic forms}, Theorem~A) therefore gives every inverse power
of the source height on these sets. Taking the largest of the
finitely many constants works on their union. At an arbitrary
point the infimum defining $\mathrm{quotHeight}$ is at most the
height of its selected rational translate, while the value of the
form is unchanged. This proves the stated quotient-height bound.
In the anisotropic case the arithmetic quotients are compact.

For $\mathrm{hc}$ intersect the prescribed open level with the
finite adelic subgroup and take a compact open subgroup inside the
intersection. At infinity the given frames identify the maximal
compact factors. The coefficient spaces of their continuous
finite-dimensional representations contain only finitely many
compact types, by complete reducibility for a compact group.
The simultaneous scalar restrictions on the Gelfand central
operators fix the infinitesimal character: for
$\mathfrak{gl}_r(\mathbb C)$ their leading symbols are the
power sums $\operatorname{tr}(X^k)$ for $1\leq k\leq r$.
These generate the invariant polynomials in characteristic zero,
so filtration by degree gives generation of the enveloping-algebra
center. For $\mathrm{fin}$ the Capelli polynomial in the source
has the elementary symmetric generators as its central coefficients
at each embedding; all of its values determine these coefficients
by interpolation. The two embeddings at a complex place give the
two factors of the complexified Lie algebra. In either case, if the
prescribed scalars cannot occur, the space is zero. Otherwise
Harish-Chandra's finiteness theorem for automorphic forms with
fixed finite level, finitely many compact types and a finite-codimension
central annihilator gives a finite-dimensional ambient space
(Harish-Chandra, \emph{Automorphic Forms on Semisimple Lie Groups};
Borel--Jacquet, \emph{Automorphic forms and automorphic
representations}, Corvallis, \S4). For $\mathrm{GL}_r$ the
split central direction is included among the central differential
operators. Its differential equation fixes the character on the
connected split center; the theorem can equally be applied on the
norm-one quotient. The span in each field is a subspace of the
corresponding finite-dimensional space.

For $\mathrm{cusp}$ take the quotient by the rational points and
the positive split center specified by $\mathrm{latticeGL}$.
The cuspidal spectral decomposition is a Hilbert direct sum of
irreducible admissible unitary representations of finite multiplicity
(Gelfand--Piatetski-Shapiro and Harish-Chandra, in the adelic
form of Borel--Jacquet, \S4). For any irreducible closed summand
take its finite-level and $K_\infty$-finite vectors. Admissibility
identifies them with its algebraic restricted tensor product of
the local smooth and Harish-Chandra modules. The archimedean
vectors are smooth and infinitesimal-character finite. Their
automorphic realizations are cuspidal smooth functions of moderate
growth, their Hilbert classes have the chosen closure, and every
vector in that summand satisfying the two finite-vector conditions
is in this subspace. Their representatives are unique by continuity
and full support of Haar measure. They are invariant under the
given split center and give precisely $\mathrm{IsFinitePartOf}$.
For the separate boundedness clause use rapid decrease of a
cuspidal automorphic form on Siegel sets modulo the split center.
The hypothesis on its absolute value allows the bound to extend
to the whole group even when its central action is nontrivial.
In rank one the idele class quotient modulo the split center is
compact. Rank zero gives the trivial group in all the relevant
assertions. This proves $\mathrm{decay,hc,fin,cusp}$.
% END A72F4

% BEGIN A72F5
\paragraph{A72F5. Isobaric uniqueness, contragredients and infinity types.}
For the finite lists in $\mathrm{js}$, equality of the multisets
at almost all places is equality of the standard unramified
Euler polynomials of their isobaric sums. Jacquet--Shalika's
isobaric uniqueness theorem (\emph{On Euler products and the
classification of automorphic forms}, I--II, 1981) gives equality
of the cuspidal supports as multisets of cuspidal isomorphism
classes, including their twists. Pair occurrences within each
class to obtain the required bijection, with equality of ranks.
This covers repetitions and empty families. For the bound,
unitarity of the central character makes the cusp unitarizable
in its cuspidal realization. Its local components are generic.
The strict Jacquet--Shalika estimate at each unramified place
gives $|\alpha|<q^{1/2}$; applying it to the unitary contragredient,
whose roots are the inverses, gives $q^{-1/2}<|\alpha|$.
For rank zero the Satake multiset is empty.

Write $\theta(g)={}^{t}g^{-1}$. The Gelfand--Kazhdan
contragredient theorem for local $\mathrm{GL}_n$ gives
$\pi_v^\theta\simeq\pi_v^\vee$ at each finite place. At a
complex place this identity also follows from the Langlands
classification, by taking the dual parameters. Normalize the
resulting evaluation pairings to take value one on the two
chosen spherical vectors at almost all places. Their product
defines a bilinear pairing on the algebraic restricted tensor
product. On any finite set of factors the product pairing is
nondegenerate on both sides; enlarging the set to contain the
factors of any given vector proves the same globally. The local
automorphic factorization transports it to the specified
$\pi.V$. Invariance gives the two finite and compact identities;
differentiating on the real group gives the sum of the two
infinitesimal actions equal to zero with second argument
$d\theta(X)=-{}^t X$. These are exactly the three action clauses
of $\mathrm{Paired}$, proving $\mathrm{gk}$, also for the
one-dimensional trivial representation in rank zero.

For $\mathrm{pur}$ the input is a cuspidal automorphic representation
whose infinitesimal exponents at every complex embedding belong
to $(n-1)/2+\mathbb Z$. Clozel's purity lemma (\emph{Motifs et
formes automorphes}, Lemma~4.9) applies to algebraic cusps,
including singular infinity types. In the shifted convention it
pairs the two multisets at conjugate embeddings by $z\mapsto w_0-z$
for a common $w_0$. Equivalently, remove the common real norm
power from the central character to make the cusp unitary.
For a unitary module the conjugate module is its contragredient;
on the two complexified Lie factors conjugation interchanges
embeddings and duality negates the shifted infinitesimal exponents.
Restoring the norm power adds the same scalar twice to each paired
sum. These sums are integers since both exponents belong to the
same half-integral coset. This proves the asserted integrality
and uniformity of $w_0$.

For distinct exponents in that coset, Clozel's cohomological
criterion for regular algebraic cusps (ibid., \S3) gives an
irreducible algebraic coefficient system with nonzero relative
Lie algebra cohomology. One can describe its infinitesimal
character by taking the dual of the algebraic highest weight
obtained after subtracting the half sum of positive roots.
Its matrix entries are regular functions generated by the
matrix coordinates at all embeddings and inverse determinants,
as required by $\mathrm{AlgRep}$. For the symmetric pair
$(\mathrm{GL}_n(\mathbb C),U(n))$ the tangent complement is
the real space of Hermitian matrices. Its brackets have zero
projection to that complement. Thus the relative differential
on equivariant alternating cochains is exactly the source's
alternating sum of coefficient actions. Tensor with a nonzero
finite-adelic vector to obtain the cochain on the full cuspidal
module. Contracting by a functional taking that vector to one
shows that a boundary there would have been a boundary on the
archimedean module. This proves $\mathrm{coh}$.
% END A72F5

% BEGIN A72F6
\paragraph{A72F6. Strong base change with the ascending Speh embedding.}
Let $F=E^+$ and take the definite unitary group of the identity
Hermitian form. Its center is anisotropic and it is compact at
infinity. Hence the given irreducible space of functions is in the
discrete spectrum. Every finite nonsplit place is unramified in
$E$; the identity lattice defines a hyperspecial subgroup there.
Moreover $[F:\mathbb Q]$ is even: the quadratic idele class
character of $E/F$ takes the value $+1$ on the local unit $-1$
at every finite place and $-1$ at every real place. Its global
value on $-1$ is one. In particular Labesse's degree restriction
is satisfied. Apply his \emph{Changement de base CM et s\'eries
discr\`etes}, Corollary~5.3, in the strong base-change formulation
of Guerberoff, \emph{Modularity lifting theorems for Galois
representations of unitary type}, Theorem~2.2. The archimedean
discrete-series and trace-isolation hypotheses there hold for
the compact real group. We get positive ranks $N_i$ with sum $n$,
and conjugate-self-dual discrete representations $\Pi_i$ of
$\mathrm{GL}_{N_i}(\mathbb A_E)$ whose isobaric sum has the
prescribed local base change at every split place and every
spherical inert place, and the transferred infinitesimal character
at infinity.

Realize each $\Pi_i$ in the discrete automorphic functions and
take its algebraic finite vectors for $V_i$. Its local components
and infinitesimal characters give $\sigma_{i,w},m_{i,\tau}$.
The discrete spectrum classification of M\oe glin--Waldspurger
(\emph{Le spectre r\'esiduel de $\mathrm{GL}(n)$}, 1989) gives
$\Pi_i=\operatorname{Speh}(\delta_i,d)$, for a unitary cusp of
rank $a$, $ad=N_i$. Uniqueness, conjugation and duality of the
cuspidal support give $\delta_i^c\simeq\delta_i^\vee$. Its central
character on the positive split center is its own inverse and is
continuous, hence trivial there. Choose its normalized cuspidal
realization $P_i$. The Speh representation embeds in the
\emph{ascending} normalized induction with exponents
$-s_i$, where $s_i=(d-1)/2-i$ for $0\leq i<d$. This is the
dual of its descending Langlands quotient realization, applied
first to the contragredient cusp. Passing to the upper block
parabolic without normalized induction adds the exponent
$a s_i$ from its modulus to block $i$. Thus the characters in
the source's $\mathrm{indForms}$ have exponents $(a-1)s_i$,
exactly $\mathrm{AscEmb}$. The intertwiner on $V_i$ can be
extended linearly to the ambient function space by a vector-space
complement, since injectivity and equivariance are tested on
$V_i$. This gives $\mathrm{CsdAsc}$ including the normalization
of $P_i$.

Take the local component of the isobaric sum for the source's
$P_w$. It is an irreducible subquotient of normalized induction
from the $\sigma_{i,w}$. At a split place the projection to the
chosen $w$ is the split identification. Local factorization with
one nonzero tensor in the other factors realizes the required
intertwiner with the original unitary space. At a spherical place
the transfer is spherical. At an inert place where a Satake tuple
$\alpha$ is given, the unramified base-change homomorphism sends
the relative parameters to the multiset consisting of all
$\alpha_j,\alpha_j^{-1}$ and one extra $1$ for odd $n$
(M\'inguez, \emph{Unramified representations of unitary groups},
\S4). The multiset is unchanged by the choices of integral
frame, Weyl representative or uniformizer in the Satake clause.
At the other inert places the subquotient already supplies all
the data required. At infinity the differential of the
algebraic representation of highest weight $\lambda$ has shifted
character $\{\lambda_k+(n-1)/2-k:0\le k<n\}$, and archimedean
base change identifies it with the union of the $m_{i,\tau}$.
This proves $\mathrm{BaseChangeIsobaricN}$ on all its inputs.
% END A72F6

% BEGIN A72F7
\paragraph{A72F7. Passing from a tested weak local correspondence to the actual parameter.}
Write $C$ for classical local Langlands with geometric reciprocity,
and $c$ for any map satisfying $\mathrm{IsRecT}$. The canonical
map satisfies these weak laws: in addition to the segment, twist,
duality and central character properties of local Langlands,
the type formula on a regular finite-field cuspidal type is the
depth-zero correspondence. In detail an irreducible depth-zero
supercuspidal is parametrized by an unramified extension and a
regular character of its residue field. On inertia its parameter
has the Frobenius conjugates of this character, while its unramified
normalization is fixed by the determinant. This is also the
degree-zero calculation used in the source's
$\mathrm{isRecT\_of\_isRec}$
(Silberger--Zink, \emph{An explicit matching theorem for
level zero discrete series of unit groups of $\mathfrak p$-adic
simple algebras}, and the depth-zero correspondence for
$\mathrm{GL}_n$). The rank-zero entry is the unique zero parameter.

Recall the two relevant uniqueness assertions. The source's
$\mathrm{recTUnique}$ identifies $c(u)=C(u)$ for good $u$:
use rank-one reciprocity on a character, the regular residue-field
type and central character on a depth-zero supercuspidal over
$\mathbb Q_p$, then the twist and segment laws on their lists.
These are exactly the allowed good supports. Separately the
injectivity axiom of $c$ says that $c(x)=c(u)$ and good $u$
imply $x=u$, with no prior condition on $x$.

Let $y=c(x)$ with zero monodromy. Under either of the guards below
take $u=C^{-1}(y)$. In the toral guard, the characteristic
polynomial of every commutator is $(T-1)^n$. A commutator belongs
to inertia, whose image in a Weil--Deligne representation is finite.
Such a unipotent element is the identity. The Frobenius-semisimple
abelian Weil representation therefore splits into characters. Its
inverse under local Langlands has support on characters and is good.
In the tame guard the field is $\mathbb Q_p$. The source tests
each wild inertia element by the same polynomial; finite order
again gives trivial action. For a tame semisimple Weil
representation the characters of tame inertia decompose into
Frobenius orbits. On one such orbit, diagonalize the power of
Frobenius which fixes one eigenspace. Each resulting simple
summand is induced from a character of the unramified extension
of degree the orbit length, with regular residue character when
that length exceeds one. Under the depth-zero correspondence
it gives a depth-zero supercuspidal up to unramified twist, and
a character in length one. Hence $u$ is good in this case also.
It follows in both cases that $c(u)=C(u)=y=c(x)$, whence $u=x$.
Thus the actual parameter equals the one tested by the arbitrary
map. The tame argument applies in particular with the squarefreeness
guard and the prime bound of the source. It gives no assertion
of unrestricted uniqueness for the weak correspondence.
% END A72F7

% BEGIN A72F8
\paragraph{A72F8. A single polarized Galois representation for all the local rows.}
Fix the original CM field, closure, rank $n\ge1$, prime, and
$\iota$. The regular infinitesimal characters give regular
algebraicity. The bilinear pairing in $\mathrm{PolarizedPair}$
identifies $\pi^c$ with $\pi^\vee\otimes(\chi\circ N_{E/F})$;
its sign clause fixes $\chi_v(-1)=(-1)^n$ at each real place.
These are a regular algebraic polarized cuspidal pair in the sense
of Barnet-Lamb--Gee--Geraghty--Taylor, \emph{Potential automorphy
and change of weight}, Theorem~2.1.1. Take its continuous
semisimple $r_{\ell,\iota}(\pi)$ once and realize it in a basis
over a finite extension of $\mathbb Q_\ell$, then extend the
coefficients. In rank one the representation and compatibilities
are given directly by algebraic Hecke reciprocity and A72F3.

Here is the weight conversion. If the ordered shifted exponents
at $\iota\tau$ are
$z_i=\lambda_i+(n-1)/2-(i-1)$ for $1\le i\le n$, then
the dual highest weight is $(-\lambda_{n+1-i})_i$. The labelled
jumps of the attachment are $i-1-\lambda_i$, equivalently
$(n-1)/2-z_i$. This uses the convention that the cyclotomic
character has jump $-1$. A72F5 supplies
$z_{\bar\tau}=w_0-z_\tau$ as multisets, so these jumps are
paired at conjugate embeddings by $h\mapsto (n-1-w_0)-h$.
This fixes a single integer weight and the whole labelled
multiset family before any place or period model is chosen.

The local normalization is $x=\pi_w\otimes|\det|^{(1-n)/2}$.
Away from $\ell$ the attachment theorem gives the
Frobenius-semisimplified local compatibility (in particular
compatibility at zero monodromy). At $w\mid\ell$ it gives de
Rhamness and the indicated labelled jumps. For monodromy use
Barnet-Lamb--Gee--Geraghty--Taylor, \emph{Local--global compatibility
for $l=p$, II}, Theorem~1.1, which applies also to essentially
conjugate-self-dual cusps. The underlying semisimple Weil modules
agree, and the monodromy partitions on the Galois side are
dominated by those for $C(x)$. When $N$ on $C(x)$ vanishes all
its blocks have size one, forcing all Galois blocks to have size
one. Thus the two Frobenius-semisimplified WD classes agree
under this hypothesis, for every regular weight. For a spherical
$\pi_w$ the parameter is unramified with $N=0$; the potentially
semistable representation is crystalline since inertia is trivial
on WD and $N=0$. Potential semistability and the labelled
calculations transport to the arbitrary supplied tied model
and bridge by A72F2.

Finally for every tested map $c$ and every supplied local component,
$\dim c(x)=n$ by its dimension law. Both Galois constructions have
dimension $n$ at the applicable places. In either guarded
zero-monodromy case A72F7 gives $c(x)=C(x)$. The WD comparison
above yields equality in the toral branch and gives zero
monodromy and equality of the characteristic polynomials on every
degree-zero Weil element in the tame squarefree branch. At
$\ell$ one also has the stipulated de Rham and potentially
semistable models and the jumps at every induced embedding.
This proves $\mathrm{PropsB}$, hence the precise $\mathrm{RowT}$
required by $\mathrm{gal}$.
% END A72F8

% BEGIN A72F9
\paragraph{A72F9. Assembly of the automorphy literature hypothesis.}
Use A72F4 for the decay, two finiteness and cuspidal-spectrum
fields. A72F5 gives Jacquet--Shalika, Gelfand--Kazhdan, purity and
cohomologicality. A72F6 supplies definite-unitary base change
and A72F8 supplies the polarized Galois row. Use A72F1 for
global Artin reciprocity and the ray field, A72F3 for the two
character comparisons and A72F2 for the de Rham implication.
These are all fifteen fields of $\mathrm{LiteratureAutomorphyHypothesis}$.
% END A72F9
